\chapter{Reference Manual}
\label{ch:geolog-reference}

\noindent\draftmark\medskip

This chapter is the systematic account of geolog as shipped in the
\path{geolog-alpha} repository: the source language and its grammar,
the normalization pipeline and its registry, the elaboration fragment
the \texttt{geolog} verb carries, the chase and its judgments, the
meta-schema, and the substrate contracts — commits, storage, sync,
the ground seam, and capabilities. The Quick Start
(chapter~\ref{ch:geolog-quickstart}) and the Tutorial
(chapter~\ref{ch:geolog-tutorial}) teach by example; this chapter
specifies. It restates none of their walkthroughs and cites them
where an example earns more than a rule would.

Three reading conventions govern everything below.

\paragraph{Claim discipline.} Normative present tense describes
behaviour shipped and tested in the repository as of this edition's
receipt passes (2026-07-21 through 2026-07-31, recorded at this
chapter's source tail). Designed-but-not-shipped behaviour appears only in
clauses explicitly marked \emph{design tense}, each naming where the
design lives today. Every quoted error message and every transcript
fragment in this chapter is the verbatim output of a real run against
the committed tree (receipts in this chapter's source tail).

\paragraph{Authority order.} The \texttt{geolog} binary is
ergonomics and composition over constituent crates, and its own crate
documentation states the rule this manual inherits: where a behaviour
differs between the facade and a constituent, the constituent is
right and the facade has a bug. Accordingly, each section below names
its authority — the crate (or committed theory file) whose documented,
tested behaviour it restates. Where this manual and that source
disagree, the source wins and the disagreement is a defect in this
manual.

\paragraph{Scope.} This manual covers the public
\path{geolog-alpha} tree. The engine library supports strictly more
than the command-line fragment (\S\ref{sec:gr-engine}); the manual
covers both, and says which surface each statement binds.

\section{Lexical structure}
\label{sec:gr-lexical}

\emph{Authority: \path{crates/geolog-normalize/src/lex.rs}.}

Theory text and instance blocks share one tokenizer. Input is
processed line by line; a \texttt{//} begins a comment (outside a
string literal) that extends to the end of the line and is
discarded. Tokens, longest match first:

\begin{Verbatim}[fontsize=\small,frame=leftline,rulecolor=\color{accentmid},framesep=3mm,xleftmargin=3mm]
two-char    ::= "->" | "|-" | "<=" | ">=" | "!="
single-char ::= "<" | ">" | "=" | "+" | "*" | "[" | "]" | "(" | ")"
              | "," | ":" | "." | "{" | "}" | ";"
integer     ::= digit+
string-lit  ::= '"' (any scalar except '"' or newline)* '"'
identifier  ::= starter continuer*
\end{Verbatim}

\noindent The identifier alphabet is Unicode-open, and the operator
characters are a closed ASCII set, so the two are disjoint by
alphabet rather than by any gating rule. A \emph{starter} is an
ASCII letter, \texttt{\_}, a non-ASCII alphabetic scalar (Greek
letters, $\ell$, \ldots), or one of the math-letterlike symbols
$\partial$ and $\nabla$; a \emph{continuer} is any starter, an ASCII
digit, \texttt{\_}, \texttt{/} (so \texttt{ax/path/edge} and
\texttt{Marking/count} are single tokens), or a script extender ---
the superscript and subscript digits and signs, and the primes.
Script extenders continue tokens and never start them: $\partial^{+}$
is one token, and a detached superscript at token start is a
tokenization refusal, so $\partial$~$^{+}$ (with a space between)
cannot silently mean anything. String literals are single-line and carry no escapes; the
grammar admits them in exactly one position --- contraction specs in
\texttt{law} bodies (\S\ref{sec:gr-editions}) --- and every other
position refuses them downstream. There are no negative-number
literals and no block comments in the theory grammar. Anything else
stops the tokenizer with a hard error naming the offending text.

The reserved words of the theory grammar are:

\begin{Verbatim}[fontsize=\small,frame=leftline,rulecolor=\color{accentmid},framesep=3mm,xleftmargin=3mm]
forall  exists  not  iff  or  false  theory  extends
Set  Prop  Index  Fin  DynTensor  List
\end{Verbatim}

\noindent \texttt{Fin} and \texttt{Index} name the same bounded-index
type (\texttt{Fin} is the ratified spelling, \texttt{Index} the
legacy alias --- \S\ref{sec:gr-decls}). \texttt{Tensor} is not
reserved, and no longer names a surface type at all: the
shape-in-type tensor former lives in the meta-schema tier
(\S\ref{sec:gr-meta-tier}), and the surface tensor type is the
keyword \texttt{DynTensor} (\S\ref{sec:gr-decls}). \texttt{instance}
is likewise not a keyword of the theory grammar: an
\texttt{instance} \emph{block} is recognized by the file splitter
(\S\ref{sec:gr-files}), one layer up, and the same word doubles as
an optional postfix marker on model-of types
(\S\ref{sec:gr-decls}), recognized contextually. The mode marks
\texttt{chase}, \texttt{report}, and \texttt{refuse}
(\S\ref{sec:gr-chase-modes}) are contextual in the same way ---
recognized only as a statement's leading token: \texttt{chase : …;}
is still an ordinary declaration named \texttt{chase}.

\section{Source files}
\label{sec:gr-files}

\emph{Authority: \path{crates/geolog/src/input.rs}.}

A \path{.geolog} file carries two strata: \emph{theory text} and
\emph{instance blocks}. Before parsing, the file is split by a
deliberately minimal scanner that tracks \texttt{//} comments and
brace depth, and recognizes the word \texttt{instance} only at brace
depth zero, at an identifier boundary on both sides. Each
\texttt{instance … \{ … \}} span is extracted whole (an unclosed
block is a hard error); the remaining text — with each extracted span
replaced by whitespace of identical line structure, so diagnostics
keep their line numbers — is handed to the normalization pipeline
verbatim. Consequences worth stating: the word \texttt{instance}
inside a comment or inside any braced body is ordinary text; theory
declarations and instance blocks may interleave freely at top level;
and the theory grammar remains entirely the normalizer's business —
the splitter never interprets theory text.

\section{Theories and declarations}
\label{sec:gr-decls}

\emph{Authority: \path{crates/geolog-normalize/src/parse.rs} (grammar)
and \path{crates/geolog-normalize/src/ast.rs} (abstract syntax).}

Theory text is a sequence of theory declarations:

\begin{Verbatim}[fontsize=\small,frame=leftline,rulecolor=\color{accentmid},framesep=3mm,xleftmargin=3mm]
theory-decl ::= "theory" NAME param* ["extends" NAME] "{" statement* "}"
param       ::= "(" NAME ":" ("Set" | NAME) ")"
statement   ::= [mode] NAME ("," NAME)* ":" statement-body ";"
mode        ::= "chase" | "report" | "refuse"
\end{Verbatim}

\noindent A theory header may bind parameters, name first: a
\texttt{(X : Set)} group binds a sort parameter, and a
\texttt{(N : PetriNet)} group --- whose second name must be a theory
declared earlier in the file --- binds a \emph{model-of} parameter,
a model of that theory. (An earlier binders-before-name header order
parses as a deprecated spelling; see \texttt{--strict},
\S\ref{sec:gr-cli}.) Every statement is an optional mode mark, one
or more comma-separated names sharing one signature, a colon, a
body, and a semicolon. The statement's \emph{kind} is dispatched on
its name: \textbf{a statement whose name begins with \texttt{ax/} is
an axiom} (\S\ref{sec:gr-axioms}); a statement opened by the
contextual keyword \texttt{law} is a law member of the aggregation
surface, whose grammar this edition defers
(\S\ref{sec:gr-editions}); every other statement is a declaration.
The prefix does grammatical work rather than expressing
a naming convention: a sequent-shaped body under a name without the
\texttt{ax/} prefix is parsed as a declaration and fails (verified:
\texttt{my/rule : forall …} stops with \texttt{syntax error:
unparseable type at line 4: forall …}, the cited line being the
statement's own). The Quick Start's remark that
slash-names are ordinary identifiers remains true — the slash is an
identifier character (\S\ref{sec:gr-lexical}) — but the two-character
prefix \texttt{ax/} is what makes a statement an axiom. The mode
mark decides whether an axiom fires or judges
(\S\ref{sec:gr-chase-modes}); the AST carries the \emph{authored}
mark, so rendering preserves the source form, and a mode mark on a
declaration head is a hard error naming both the mark and the head
(the message names \texttt{ax/} and \texttt{law} as the two homes a
mode can mark).

Declaration bodies:

\begin{Verbatim}[fontsize=\small,frame=leftline,rulecolor=\color{accentmid},framesep=3mm,xleftmargin=3mm]
sort-decl   ::= NAME ":" "Set" ";"
record-decl ::= NAME ":" "{" field ("," field)* "}" "->" codomain ";"
unary-decl  ::= NAME ":" type "->" codomain ";"
curried-decl::= NAME ":" ("(" NAME ":" type ")" "->")+ codomain ";"
const-decl  ::= NAME ":" type ";"
field       ::= NAME ":" type
codomain    ::= "Prop" | type
type        ::= NAME                     -- a sort (or a theory: model-of)
              | term ("Fin" | "Index")   -- bounded index (n elements)
              | NAME "DynTensor"         -- dynamic tensor over an element sort
              | NAME "List"              -- finite sequences over an element sort
              | arg+ NAME ["instance"]   -- model-of application
\end{Verbatim}

\noindent Braces name record domains. (The bracket spelling
\texttt{[f:\,v]} of earlier corpora still parses by default while
the committed corpus ports to braces; under strict parsing it is a
deprecation error pointing at the brace form, and \texttt{geolog
fmt} prints the canonical respelling --- \S\ref{sec:gr-cli}.) A
curried signature of named binders normalizes to the record form at
parse, so the two spellings declare the same relation or function;
an unnamed curried binder is a hard error. A declaration with
codomain \texttt{Prop} declares a
\emph{relation}; any other codomain declares an element-valued
\emph{function} (a \texttt{const-decl} is the nullary case: a
constant of its sort). Relations come in two domain shapes at the
surface: \emph{record} (named fields; field order as written is
preserved for display, but fields are name-keyed, never positional)
and \emph{unary} (one anonymous argument). There is deliberately no
surface form for a zero-argument proposition: \texttt{Flag : Prop;}
parses as a constant declaration of a (nonexistent) sort
\texttt{Prop} and is refused downstream (\S\ref{sec:gr-elab}). A
proposition wants at least one argument, or it wants to be a fact.

Two type forms are new with the aggregation surface, and both are
codomain vocabulary for the loader's stratum rather than the chase
fragment's. \texttt{DynTensor} is the \emph{dynamic} tensor type
over an element sort — shape carried in the value, where the
meta-tier \texttt{Tensor} former carries shape in the type
(\S\ref{sec:gr-meta-tier}); tensor values are side-condition
payload, never term-language citizens. A theory name in type
position means ``a model of that theory'': bare (\texttt{base :
G;} under a \texttt{(G : Graph)} parameter), or applied to
arguments (\texttt{initial\_marking : n0 Marking;}), with an
optional postfix \texttt{instance} keyword as an explicitness
marker — neither spelling is canonical, and the source printer
preserves the author's choice. A third form belongs to
the list former (\S\ref{sec:gr-meta-tier}): \texttt{Str List} is the
finite-sequence type over an element sort, postfix like
\texttt{Index} and \texttt{DynTensor} and taking one element sort
only (\texttt{Str List List} is a syntax error), with list values
consumed through the \texttt{len}/\texttt{at}/\texttt{concat} ground
roster as side-condition payload — never term-language citizens. The
surface element scope is the builtin ground sorts: a list codomain
at a declared sort is refused by the normalizer (verified:
\texttt{List element "A" is not a builtin ground sort (v1 surface
scope)}), and the wire tier's wider element scope, nesting included,
arrives with its own surface spelling rather than by fall-through.
The v1 chase bridge refuses all three forms (\S\ref{sec:gr-elab}).

\texttt{extends} names a parent theory. In the single-file pipeline
the parent must be declared in the same file, and extension is
realized by elaborating the whole file as one union
(\S\ref{sec:gr-elab}); in the committed-revision pipeline extension
is a \texttt{TheoryExtends} row with its own conformance rules
(\S\ref{sec:gr-meta}).

\paragraph{Builtin ground sorts.} The authoring idiom accepts a
roster of ground sorts in declarations without a \texttt{Set}
statement — \texttt{Q}, \texttt{Z}, \texttt{Nat}, \texttt{B},
\texttt{Str}, \texttt{ENNRat}, \texttt{Shape}, \texttt{Timestamp},
\texttt{Bytes}, \texttt{PeerId}, \texttt{CommitId}, and the
aggregation surface's staged registry sort \texttt{Accum}
(\path{crates/geolog-normalize/src/builtin.rs}) — together with the
literal-dtype correspondence that types bare integer literals
(\texttt{q}/\texttt{z}/\texttt{nat}/…), the builtin comparison
and arithmetic vocabulary the grounding moves compile onto
(\S\ref{sec:gr-normalize}), and the aggregation tier's staged
arithmetic (\texttt{nat\_mul}, \texttt{succ}). This roster is,
\texttt{Accum} aside, the normalizer's \emph{accepted} subset of the
engine's primordial sort tier (\S\ref{sec:gr-meta-tier}) —
\texttt{Accum} belongs to the deferred aggregation surface
(\S\ref{sec:gr-editions}), not the primordial tier; the command-line
chase fragment is narrower still and refuses ground sorts at
elaboration (\S\ref{sec:gr-elab}).

\section{Axioms and geometric normal form}
\label{sec:gr-axioms}

\emph{Authority: \path{crates/geolog-normalize} (\path{parse.rs},
\path{normalize.rs}, \path{lint.rs}); the normal form is also stated,
as schema, by the meta-schema (\S\ref{sec:gr-meta}).}

\subsection{The axiom grammar}
\label{sec:gr-axiom-grammar}

\begin{Verbatim}[fontsize=\small,frame=leftline,rulecolor=\color{accentmid},framesep=3mm,xleftmargin=3mm]
axiom      ::= [mode] AXNAME ":" ["forall" binders "."] [premise] "|-" conclusion ";"
binders    ::= NAME ":" type ("," NAME ":" type)*
premise    ::= item ("," item)*
item       ::= atom
             | "not" atom
             | "(" item ")"
             | "exists" binders "." premise          -- premise only
             | "forall" binders "." iff-group        -- pattern moves
conclusion ::= "false"
             | disjunct ("or" disjunct)*
disjunct   ::= ["exists" binders "."] item ("," item)*
atom       ::= term [cmpop term]
cmpop      ::= "=" | "!=" | "<" | "<=" | ">" | ">="
term       ::= prod ("+" prod)*
prod       ::= postfix ("*" postfix)*
postfix    ::= primary NAME*                          -- application chain
primary    ::= "{" NAME ":" term ("," NAME ":" term)* "}"   -- record
             | INTEGER                                       -- literal
             | NAME                                          -- variable
\end{Verbatim}

\noindent Grammar facts with semantic weight: \texttt{or} binds
loosest in a conclusion, and an existential's scope never spans an
\texttt{or} — each disjunct carries its own \texttt{exists} prefix,
so the surface conclusion is literally the normal form's shape (a
disjunction \emph{of} existentially quantified conjunctions). An
\texttt{exists} group in a premise swallows the remaining premise
items up to the turnstile. Terms are postfix application chains —
\texttt{m p Marking/count} reads ``\texttt{Marking/count} applied to
(\texttt{p} applied to \texttt{m})'' — with \texttt{*} binding
tighter than \texttt{+}, record literals in the same brace idiom
declarations use (brackets remain the parsed legacy spelling during
the corpus port, \S\ref{sec:gr-decls}; on the aggregation surface
square brackets instead delimit contraction index groups in
\texttt{law} bodies, \S\ref{sec:gr-editions}), and bare integer
literals (typed later by the
normalizer). \texttt{|- false} declares a \emph{denial}: a rejection
rule, the empty disjunction.

\subsection{Geometric normal form}
\label{sec:gr-gnf}

The engine consumes \emph{geometric} theories, and the normal form is
the central definition of this chapter:

\begin{definition}[geometric sequent]
\label{def:gr-geometric}
A geometric sequent is a universally quantified entailment whose
premise is a finite conjunction of atoms and whose conclusion is a
finite disjunction of existentially quantified finite conjunctions of
atoms; the empty disjunction is $\bot$ (a denial). An atom is an
equation between terms or a relation applied to a term; the
quantified variables are typed at declared sorts.
\end{definition}

\noindent Geometric sequents are the sequent fragment of geometric
logic \cite[Part~D1]{johnstone2002}; database theory knows the
single-disjunct case as tuple- and equality-generating dependencies
\cite{fagin2005data}. The authoring idiom of
\S\ref{sec:gr-axiom-grammar} is richer than this — \texttt{not}, \texttt{exists} in premises,
\texttt{forall}-\texttt{iff} groups, comparison operators, arithmetic
sugar, \texttt{Index} types. The normalizer
(\S\ref{sec:gr-normalize}) compiles every departure down to
Definition~\ref{def:gr-geometric} by a registered move or rejects it
loudly; the geometric linter then mechanically enforces the normal
form on the normalizer's own output — premise a conjunction of atoms;
conclusion a disjunction of existentially quantified conjunctions; no
residual \texttt{Index}; no untyped literal; every relation declared
\emph{and} classified (\S\ref{sec:gr-registry}); every atom
well-sorted. Pipeline failures are collected and reported together,
never truncated to the first.

\subsection{The normalization moves}
\label{sec:gr-normalize}

Each move is either total on its registered pattern or a loud error;
nothing passes through half-rewritten. In source order of the
pipeline:

\begin{itemize}
\item \textbf{Premise negation}, discharged by atom class. A negated
\emph{side-condition} atom needs no logic at all: side-conditions
are decision procedures, and a decision procedure is
negation-complete, its complement being the same evaluator with the
verdict flipped. The discharge is therefore exact and in place ---
$\lnot S \land P \vdash C$ becomes $\mathit{not/S} \land P \vdash C$,
with the complement relation registered automatically (move census:
\texttt{negated premise side-condition S discharges in place as
not/S}). A negated equation at a ground comparison sort has its
complement in the builtin vocabulary already: $\lnot(a = b)$ at
\texttt{Q} or \texttt{Nat} becomes \texttt{q\_ne} or \texttt{nat\_ne},
in place. The one genuinely logical case is a negated \emph{ordinary}
atom, or an equation at a declared sort. Ordinary vocabulary is
derivation-defined, so no procedure can decide its complement;
instead the atom \emph{crosses the turnstile} as a new leading
conclusion disjunct. $\lnot R \land P \vdash C$ becomes
$P \vdash R \lor C$: the classical-to-geometric compilation (move
census: \texttt{negated premise ordinary atom R crosses the
turnstile}). Read at the fixpoint, the crossed sequent says exactly
what the classical sentence said: wherever $P$ holds, $R$ holds or
$C$ does. What it deliberately does not offer is
negation-as-failure. The crossed head is a disjunction, so the axiom
is no longer chase-eligible (\S\ref{sec:gr-chase-classify}), and a
\texttt{chase} mark on it is a load refusal: a rule that fired on the
\emph{current} absence of a derivable row would be unsound against
every later derivation, so no move is owed to that reading. Under
\texttt{report}, the crossed disjunct is judged like any other menu
(\S\ref{sec:gr-chase-judgments}).
\item \textbf{Pattern rewrites} (\texttt{forall}-\texttt{iff}
groups). Extensionality, difference, and intersection shapes over a
registered membership relation rewrite to auxiliary mask
side-condition atoms; \texttt{not} on a membership atom in a
conclusion rewrites to a registered negation-free witness atom. All
four require registry entries (\S\ref{sec:gr-registry}); an
unregistered shape is a hard error, because auxiliary vocabulary is
an obligation on downstream evaluators and is never minted silently.
\item \textbf{Exists-hoist.} An \texttt{exists} group in a premise
hoists to universal variables — sequent semantics:
$(\exists x.\,P) \vdash C \;\equiv\; \forall x.\,(P \vdash C)$.
\item \textbf{Index erasure.} Every \texttt{n Index} type erases to
\texttt{Nat} plus a generated range obligation (a binder guard atom,
or a generated \texttt{ax/range/*} axiom).
\item \textbf{Grounding.} Comparison and arithmetic sugar compiles
onto the builtin roster at ground sorts: \texttt{!=} to
\texttt{q\_ne}/\texttt{nat\_ne}, \texttt{<}/\texttt{<=} to
\texttt{q\_lt}/\texttt{q\_le}/\texttt{nat\_lt}, \texttt{+}/\texttt{*}
to \texttt{q\_add}/\texttt{q\_mul}, with \texttt{>}/\texttt{>=}
flipped by operand swap. There is no grounding at declared sorts:
\texttt{a != b} over a declared sort \texttt{A} is a hard error
(verified: \texttt{no ground builtin for != at sort A (the !=@A case
would need a builtin-roster addition)}). Disequality of declared-sort
elements is not expressible in the fragment — deliberately: the
chase's equality discipline is congruence
(\S\ref{sec:gr-chase-congruence}), and a general disequality atom
over invented elements is exactly the kind of vocabulary that would
need an evaluator nobody has priced. (A negated premise
\emph{equation} at a declared sort is a different matter: it crosses
the turnstile as an equality conclusion, per the premise-negation
move above, so the classical constraint remains sayable even though
a disequality atom does not exist.)
\item \textbf{Literal typing.} Bare integer literals receive their
dtype from context (\texttt{q:3/1}, \texttt{nat:2}, …); an
untypeable literal fails the linter.
\end{itemize}

\noindent An axiom touched by any move is renamed with an
\texttt{\_n} suffix, so normalized output re-normalizes to itself
with an empty move log. The pipeline's guarantees, carried by its
test suite: two-run determinism (byte-identical listings),
idempotence (normalize $\fatsemi$ normalize $=$ normalize),
linter-cleanliness of every output, ground meaning preservation (the
grounding moves agree with direct evaluation of the source operators
on random ground instances), and committed fixture goldens per move.
\texttt{iff} at sequent top level (which would become two sequents)
is unimplemented and errors loudly if it ever fires.

\subsection{Imports, classification, and roster totality}
\label{sec:gr-registry}

\emph{Authority: \path{crates/geolog-normalize/src/imports.rs}
(the import surface), \path{crates/geolog-normalize/src/registry.rs}
(the registry value both surfaces feed).}

Every declared \texttt{Prop}-valued relation is classified, exactly
one of:

\begin{itemize}
\item \texttt{ordinary} — derivation-defined: populated by assertion
and by the chase, and by nothing else;
\item \texttt{side-condition} — decided by a bounded external
procedure (the ground seam, \S\ref{sec:gr-ground}, or a registered
auxiliary evaluator), never derivation-defined.
\end{itemize}

\noindent This is \emph{roster totality}: every relation has a
decided evaluator, and it holds \emph{by construction} — the
obligation is discharged by syntax instead of by hand. The
classification lives in the source file itself. \texttt{ordinary} is
the default, declared by writing nothing; a side-condition is
declared by an \texttt{import} item naming its provider:

\begin{Verbatim}[fontsize=\small,frame=leftline,rulecolor=\color{accentmid},framesep=3mm,xleftmargin=3mm]
import std.{Q, q_lt};
import <provider>.{<relation>, ...};
\end{Verbatim}

\noindent The two namespaces do different work. An
\texttt{import std} item authorizes \emph{builtin ground vocabulary}
used by name — sorts like \texttt{Q} and \texttt{Nat}, comparisons
like \texttt{q\_lt} — and a name outside the builtin roster refuses
at once, listing the importable names. Any other namespace declares
each named relation a side-condition awaiting that provider, bound
at the ground seam (an unknown provider is not an error at check
time: binding is deferred, and an unbound side-condition atom simply
cannot be discharged — refusal over guessing, downstream). Import
items are position-free top-level items, so multi-file closures
concatenate; re-importing a name under the same namespace is
idempotent, and importing one name under two namespaces refuses ---
one name, one provider. A relation or function declaration that
reuses a builtin \emph{vocabulary} name is legal only as an
exact-signature shim of that builtin: same-named, byte-equivalent to
what the engine itself would say. A builtin \emph{sort's} name admits
no shim at all — it is fully reserved, and a local sort takes another
name (verified: \texttt{sort Q (declared in T): builtin ground sort
names are reserved — the builtin Q is ambient (import std.\{Q\};
declares the usage), so a local sort must take another name}). Both
reservations are kind-local (the sort and vocabulary namespaces never
meet), so a relation may take a builtin sort's name, and a declared
sort may take a builtin relation's. Builtin
comparison relations and the normalizer's own minted auxiliaries are
side-conditions without import lines. Under \texttt{--strict},
builtin ground vocabulary used by name demands its \texttt{import
std} header even in files that otherwise import nothing.

The \emph{pattern registry} is the second classification surface,
today a library-level one: it names, per rewritten relation, its
\emph{membership} reading (which field is subject, which is index),
and the auxiliary relation each registered move targets
(\texttt{extensionality}: 2 fields; \texttt{difference} and
\texttt{intersection}: 3; \texttt{negation-witness}: the relation's
own arity). Registry construction validates all of this — missing
membership entries, unregistered auxiliaries, arity mismatches, and
duplicate auxiliaries are errors before any theory text is touched.
Pattern registrations have no source-level surface yet; the
library's registry value carries them (the normalizer's
\path{fixtures/} include pattern-bearing specimens).

A line-oriented registry \emph{file} (\texttt{--registry}) was the
classification's original carrier; it is deprecated, retained for
the migration era only. The one-authority law binds at the
registry's content, and classification is the content it guards: a
registry whose lines classify relations refuses to combine with a
source that declares imports (verified: \texttt{this source
declares imports AND an external registry carries relation
classification — classification has one home; …}), while a
pattern-only registry file composes with an importing source — its
aux names kept clear of the imported ones (verified: \texttt{aux
relation … collides with an imported name — aux mask vocabulary is
the normalizer's, never imported}), a collision being
classification smuggled back in aux costume. The
line format also carries the pattern registrations of the paragraph
above (their one textual carrier while patterns have no in-source
spelling), and a pattern registration is no second classification
authority.

\section{Instance blocks}
\label{sec:gr-instances}

\emph{Authority: \path{crates/geolog/src/facts.rs} (grammar),
\path{crates/geolog/src/seed.rs} (semantics).}

\begin{Verbatim}[fontsize=\small,frame=leftline,rulecolor=\color{accentmid},framesep=3mm,xleftmargin=3mm]
instance-block ::= "instance" NAME ":" NAME "{" item* "}"
item           ::= elem-decl | record-fact | bare-fact
elem-decl      ::= NAME ("," NAME)* ":" NAME ";"
record-fact    ::= "{" NAME ":" NAME ("," NAME ":" NAME)* "}" NAME ";"
bare-fact      ::= NAME NAME ";"
\end{Verbatim}

\noindent Tokenization is the theory lexer's, so comments and
identifiers mean the same thing in both strata. The header names the
block and the theory it instantiates; the named theory must be
declared in the same file. Items are element declarations and facts;
a record fact's brace lists \texttt{field: element} pairs and then
the relation; a bare fact (\texttt{a Marked;}) applies a unary-kind
relation. Facts must be written in the relation's own shape — a
record fact on a unary relation, or vice versa, is an error naming
the mismatch. Braces are canonical here as in declarations: this
grammar is what \texttt{geolog fmt} prints, and the formatter's
output re-parses (verified: the respelled Quick Start file chases
end-to-end, its report byte-identical to the bracket original's),
while the bracket spelling \texttt{[f:\,a]\,Rel;} remains the
accepted legacy form until the committed corpus ports to braces.

Semantics, in seeding order (source order — element numbering in
transcripts is deterministic):

\begin{itemize}
\item \textbf{Elements are declared before use.} A fact referencing
an undeclared name is an error, never a silently minted element
(verified: \texttt{`P` fact references `ghost`, which is not a
declared element — declare it first (`ghost : A;`); elements are
never minted silently}).
\item \textbf{Fields are exact.} A record fact must supply every
declared field, no extras, no repeats, no defaults; each element's
sort must equal the field's declared sort.
\item \textbf{Union semantics across blocks.} An element name
declared in two blocks at the same sort denotes one element;
duplicate facts are idempotent. A name redeclared at a
\emph{different} sort is an error (verified: \texttt{element `a`
redeclared at sort `C` but already declared at sort `A`}). This is
the merge story — set union of commits, then re-derivation
(\S\ref{sec:gr-sync-merge}) — at authoring scale; the Tutorial's
two-writer example is built on it.
\end{itemize}

\noindent Surface syntax is a view: this block form is the command
line's authoring view of an instance, and the durable form is signed
commits with content-addressed element identity
(\S\ref{sec:gr-commits}).

\section{Elaboration: the v1 command-line fragment}
\label{sec:gr-elab}

\emph{Authority: \path{crates/geolog/src/elab.rs}; the crate-level
docs of \path{crates/geolog/src/lib.rs} state the fragment.}

The bridge from the normalized authoring idiom to the engine IR
covers, in v1: \textbf{sorts}, \textbf{\texttt{Prop}-valued
relations} (record or unary domains), and \textbf{geometric axioms
whose atoms apply relations to variables}, equality conclusions
included. Everything else the idiom can express is refused with an
error naming the construct and, where one exists, the supported
spelling of the same intent. The engine library supports more than
this bridge; the bridge grows by deliberate widening, never by silent
acceptance.

Elaboration takes the file's theory declarations as \emph{one
single-file closure}: the union of every declaration in the file,
with \texttt{extends} between in-file theories honoured by that same
union. Consequences: an instance block may reference any declaration
in the file's closure regardless of which theory block carries it;
duplicate sort or relation names anywhere in the closure are errors;
and \texttt{extends} naming an out-of-file parent is refused
(verified: \texttt{theory `Child` extends `Ghost`, which is not
declared in this file — cross-revision extension rides the
committed-revision pipeline (geolog-meta-rows / geolog-revision), not
the v1 CLI bridge}).

The operative refusals, each verified against the committed binary:

\begin{itemize}
\item \textbf{Element-valued functions} (any declaration with a
non-\texttt{Prop} codomain, constants included):
\begin{Verbatim}[fontsize=\small,frame=leftline,rulecolor=\color{accentmid},framesep=3mm,xleftmargin=3mm]
function symbol `f` (codomain `C`): element-valued functions are engine
IR the v1 bridge does not carry — spell the function as a relation plus
a functionality axiom (two rows in the premise, an equality in the
conclusion), the idiom the committed theories use (e.g. `Marking/count`
+ `ax/count/functional` in theories/PetriNet.geolog)
\end{Verbatim}
The named idiom is normative: a function $f : A \to B$ is spelled as
a relation $F[a{:}\,A, b{:}\,B]$ plus
$\forall a, b_1, b_2.\; F[a, b_1], F[a, b_2] \vdash b_1 = b_2$, and
the chase's congruence closure does the rest
(\S\ref{sec:gr-chase-congruence}). \texttt{Flag : Prop;} lands in
this refusal too (its ``codomain'' parses as the sort
\texttt{Prop}); see \S\ref{sec:gr-decls} for why no zero-argument
proposition form exists.
\item \textbf{Ground sorts in declarations} — any field, argument, or
binder sort not declared \texttt{Set} in the file, the builtin roster
included:
\begin{Verbatim}[fontsize=\small,frame=leftline,rulecolor=\color{accentmid},framesep=3mm,xleftmargin=3mm]
relation `R` field `n` names sort `Nat`, which this file does not
declare — the v1 bridge elaborates single-file closures over declared
`Set` sorts only (ground sorts such as `Nat` arrive with the ground
seam; see crates/geolog-ground)
\end{Verbatim}
This is also where \texttt{Index} types land operationally: the
normalizer erases \texttt{n Index} to \texttt{Nat}
(\S\ref{sec:gr-normalize}), and the bridge then refuses \texttt{Nat}
as above. The worked discipline for doing real arithmetic
\emph{without} the seam — counts as explicit elements on a ladder,
chase-derived — is \path{theories/PetriNet.geolog}.
\item \textbf{Tensor, list, and model-of types} — the codomain
forms of \S\ref{sec:gr-decls} beyond plain sorts are loader
vocabulary, and the bridge carries a dedicated refusal for each. A
\texttt{DynTensor} anywhere in a declaration:
\begin{Verbatim}[fontsize=\small,frame=leftline,rulecolor=\color{accentmid},framesep=3mm,xleftmargin=3mm]
relation `P` field `x` uses a `Tensor` type — tensor payloads are
side-condition territory, not v1 bridge vocabulary
\end{Verbatim}
A model-of member, on today's routes, is caught upstream by the
element-valued-function refusal instead (its codomain is not
\texttt{Prop}); the bridge's own instance-valued arm — refusing
instance members as loader territory — stands behind it as
defensive depth, per the posture above. A \texttt{List} codomain is
caught upstream the same way; a \texttt{List} in a \emph{field}
position reaches the bridge's own list arm (verified:
\texttt{relation `P` field `x` uses a `List` type — list payloads
are side-condition territory (the \{len, at, concat\} roster), not
v1 bridge vocabulary}).
\item \textbf{Ground literals and term-level application in axiom
atoms}, and nested records in argument positions: refused naming the
axiom. Under the current pipeline these arms are defensive depth —
declaration-level refusals and normalizer typing catch the practical
routes first — and they stay, per the bridge's stated posture: a
bridge that assumes its input is well-formed fails silently the day
the assumption breaks.
\end{itemize}

\noindent Axiom-side checks (variables bound by \texttt{forall} or an
enclosing \texttt{exists}; sorts agreeing with declarations; fields
exact; no variable bound twice; no existential shadowing an outer
variable) are belt-and-braces re-verification of what the normalizer
already established, kept for the same reason.

\section{The chase}
\label{sec:gr-chase}

\emph{Authority: \path{crates/geolog-chase} — \path{classify.rs},
\path{chase.rs}, \path{instance.rs}, \path{judgment.rs}. The engine
has no silent path: every outcome below is either a derived row, a
typed refusal, or a reported judgment.}

\subsection{Sequent classification}
\label{sec:gr-chase-classify}

Repair arity is computed from syntax alone, by one total function,
and every engine consults it — no runtime data participates. The
conclusion is normalized to disjunctive head form ($\mathrm{Conj}$
distributes over $\mathrm{Disj}$; a $\bot$ conjunct annihilates its
branch), and:

\begin{center}
\begin{tabular}{lll}
\toprule
head disjuncts & class & chase posture \\
\midrule
$0$ & denial & never fires; violations are judgments \\
$1$ & canonical & fires (invents witnesses) — the regular fragment \\
$\geq 2$ & repair menu & never auto-fired; surfaced as a branch point \\
\bottomrule
\end{tabular}
\end{center}

\noindent One carve-out is permanent: a head that reaches an
existential over an \emph{immutable} sort (an infinitary disjunction
— the engine IR's \texttt{InfDisj}) is never chase-eligible
regardless of arity, because an infinitary disjunction is in no way a
unary disjunction; no witness-invention strategy for it can ever be
sound. The refusal is loud and names the axiom.

\subsection{Modes: the author's word on firing}
\label{sec:gr-chase-modes}

Whether an axiom fires is not inferred from its shape: it is
\emph{authored}, per axiom, by the mode mark of
\S\ref{sec:gr-decls}. The three marks:

\begin{itemize}
\item \texttt{chase} — a derivation rule: the axiom fires, inventing
witnesses and asserting conclusions. Only a chase-eligible axiom
(\S\ref{sec:gr-chase-classify}: canonical head, no \texttt{InfDisj})
may carry the mark; \texttt{chase} on an ineligible axiom is a load
refusal naming the axiom by its \emph{source} name (verified:
\texttt{axiom ax/pick/one is marked `chase` but is not chase-eligible
…; mark it `report` (judge) or `refuse` (gate), or reshape its
head}).
\item \texttt{report} — a constraint: the axiom never fires; its
violations and would-be firings surface as judgments
(\S\ref{sec:gr-chase-judgments}).
\item \texttt{refuse} — a constraint that additionally gates local
authoring: \texttt{report}'s judgment semantics plus the local write
gate of \S\ref{sec:gr-engine}. On the file path there is no commit to
gate, so its violations render exactly as \texttt{report}'s.
\end{itemize}

\noindent \textbf{An unmarked axiom means \texttt{report}} — the
least-dangerous default (\texttt{chase} creates rows and
\texttt{refuse} blocks commits; \texttt{report} only judges), so
derivation is always an explicit authorial act. Denials sit naturally
under \texttt{report}: a denial never fires in any mode, and its
violations are phase-2 judgments (\S\ref{sec:gr-chase-run})
regardless.

The marks map onto the engine's per-axiom trichotomy
(\S\ref{sec:gr-engine}): \texttt{chase} $\to$ Chased, \texttt{report}
$\to$ Monitored, \texttt{refuse} $\to$ Enforced. \texttt{geolog
chase} runs a file under this authored partition by default,
recovering each axiom's mark from the normalized declarations with
three alignment checks per axiom (count, name, binder vector) so any
drift between the normalizer's and the bridge's axiom order is a
loud refusal, never a silent mode reassignment.

\paragraph{The legacy flattener.} \texttt{--mode chased|monitored}
(\S\ref{sec:gr-cli}) survives for corpora not yet carrying explicit
marks, applying one run-level policy to every axiom: under
\texttt{chased}, every eligible axiom fires and a disjunctive head is
a theory-load error before the loop ever runs (verified, on an
unmarked file: \texttt{theory load refused: axiom ax\#0 is not
Chased-eligible: repair arity is Menu(2); only a single-disjunct
(canonical) head with no InfDisj binder may run under mode Chased —
declare it Monitored instead, or split/simplify the axiom}); under
\texttt{monitored}, eligible axioms fire and disjunctive heads
surface as branch points. On a file carrying \emph{any} explicit mark
the flag refuses, naming the marked axioms (verified: \texttt{--mode
conflicts with this file's explicit surface marks: ax/path/edge
(chase), … — drop the flag to honor the marks}): a mark is the
author's word, and a global flag that silently promoted a
\texttt{report}-marked axiom into the firing set would be the exact
hazard the marks exist to close — worst case, an equality conclusion
merging the very elements it was written to police. Where the flag
does run, the flattening is announced on stderr.

\noindent Engine-level load refusals name axioms positionally
(\texttt{ax\#i}): engine sequents carry no name field — names are
metadata, held by the caller (or, in the durable pipeline, by
meta-schema rows) — and the facade maps indices back to source names
wherever it holds them (\S\ref{sec:gr-report}); the bare positional
form above is the legacy flattener's path showing the engine's own
words.

\subsection{The two-phase run}
\label{sec:gr-chase-run}

A chase run has exactly two phases, and the separation exists for
one reason: \emph{a transient judgment can never be persisted as
if final.}

\textbf{Phase 1} fires only canonical axioms, to a fixpoint or the
iteration bound. The loop's continuation condition is instance
mutation only — discovering a violation is never a reason to keep (or
stop) iterating. \textbf{Phase 2} is one evaluation pass of every
denial and every menu axiom against the \emph{final} canonical
instance, bindings resolved through the final congruence. A menu
whose missing witness is supplied by the very last canonical firing
is therefore never recorded stale, and every judgment binding is a
canonical representative of the final fixpoint.

Hitting the bound is not an error: the instance-so-far is returned, a
non-termination judgment is recorded (with both the iterations run
and the bound), and phase 2 still reports whatever holds of the
partial instance — validity is explicitly indeterminate, not omitted
(verified: \texttt{no fixpoint within budget: 5 iterations (bound 5)
— raise --budget, or the theory does not saturate on this instance
(unbounded generation is a real outcome, reported rather than guessed
at)}).

\subsection{Restricted semantics and invention}
\label{sec:gr-chase-restricted}

Canonical firing is the \emph{restricted} chase
\cite{deutschnashremmel2008}: for each binding of
an axiom's universal context that satisfies the premise, the engine
first evaluates the conclusion as a query under that binding —
\textbf{if any witness already exists, it does not fire}. Only a
witnessless binding invents: one fresh element per existential
binder, then the conclusion's atoms are asserted (relation rows as
derived rows; equations via congruence,
\S\ref{sec:gr-chase-congruence}). The Tutorial's shipment example is
the behaviour in one line: the order with a recorded manifest got no
invention; the order without one got exactly one fresh
\texttt{?Shipment}.

Every element carries an immutable creation-provenance record,
written once and never rewritten by any merge: seeded (EDB) elements
carry caller-supplied provenance bytes (at command-line scale, the
declared name; in the durable pipeline, a commit-addressed leaf), and
chase-invented elements carry the firing axiom's index, the
context-bound elements of the firing (in context order), and the
fresh binder's ordinal within the axiom. Derived rows are
distinguished from asserted rows throughout
(\S\ref{sec:gr-report}). One invariant worth its ink: for a premise
whose own existential (hoisted per \S\ref{sec:gr-normalize}) admits
several matches over the same context binding, the first firing's
witness satisfies the rest — so exactly one witness, with a
context-only provenance term, ever exists per context binding.

\subsection{Equality and congruence}
\label{sec:gr-chase-congruence}

Equality is computational, not annotational: the instance maintains a
union-find congruence closure in the style of Nieuwenhuis and
Oliveras \cite{nieuwenhuisoliveras2007}. An
equational conclusion firing merges two elements' classes; everything
known about either — rows, provenance-bearing derivations, later
firings — holds of the merged class. Tuples are stored raw and
canonicalized through the current union-find at every read, so a
merge never rewrites tables. The report renders a merged class by
joining its seeded names (\texttt{sf\_record=wh\_label};
\texttt{b=c}) and accounts for it in the carrier arithmetic
(\texttt{3 seeded, 1 merged by congruence} — verified transcript).
The functionality idiom of \S\ref{sec:gr-elab} plus this machinery is
how geolog has functions without function symbols in the v1
fragment. (Chase steps interleaved with congruence closure are also
the computational core of datalog-with-equality and
equality-saturation engines \cite{bidlingmaier23,zhang23}, the
nearest kin to this engine in the evaluation literature.)

\subsection{Judgments}
\label{sec:gr-chase-judgments}

A judgment is one thing the chase found that a silent-success engine
would have swallowed — the authority module's own one-line definition
of the type. The chase's loud outputs, exhaustively, each carrying
the axiom's identity and, where applicable, the witnessing binding
with canonical elements:

\begin{itemize}
\item \textbf{Denial violation} — a denial's premise is satisfied at
the fixpoint: $\vdash \bot$ for this binding. Reported, never
dropped, never repaired.
\item \textbf{Branch point} — a menu axiom's premise is satisfied
with no existing witness for \emph{any} disjunct. The payload carries
the menu (the disjuncts themselves, as data); the engine picks
nothing. Under the legacy run-level flattener a menu has length
$\geq 2$; menu-of-one branch points (a forced repair, surfaced rather
than fired — \texttt{a menu of 1 repair}) arise wherever the
per-axiom partition holds an \emph{eligible} axiom to \texttt{report}
(\S\ref{sec:gr-chase-modes}, \S\ref{sec:gr-engine}) — an unmarked
derivation rule's would-be firings surface this way.
\item \textbf{Non-termination} — the bound was hit
(\S\ref{sec:gr-chase-run}).
\end{itemize}

\noindent Judgments are results, not errors: a run that produces
them completed, exits with the judgment exit code
(\S\ref{sec:gr-cli}), and its instance is real — the event-sourcing
posture (invalid states representable, detected, repaired by further
ordinary commits) held formally. The Tutorial's merge chapter is the
worked argument for why this posture, and not write-rejection or
automatic repair, is the design.

\subsection{The report}
\label{sec:gr-report}

\emph{Authority: \path{crates/geolog/src/report.rs}.}

The \texttt{chase} verb renders every carrier, every row, and every
judgment, through the source file's own names; nothing about a
judgment is summarized away. The normative rendering facts, seven of
them:

\begin{itemize}
\item each sort prints its canonical carrier size with the seeded
count, the count invented by the chase, and the count merged by
congruence, followed by a \texttt{merged:} line naming merged
seeded-name groups;
\item each relation prints in declaration order, rows sorted and
deduplicated after canonicalization, with seeded/derived tallies;
\item derived rows — rows not in the canonicalized image of the
seeded baseline — are marked \texttt{+};
\item rows print in the relation's own shape (record, bare, or bare
name for nullary);
\item chase-invented elements render as \texttt{?Sort.n} in
first-encounter order;
\item side-condition relations are tagged \texttt{[side-condition]};
\item judgments print the axiom's \emph{source} name and the binding
in context order, a branch point adds its menu size, and a clean run
ends \texttt{fixpoint: reached, no judgements — the model is % spell-ok: judgement (geolog's report.rs output, quoted verbatim)
saturated.} (with a count of invented elements when there are any).
\end{itemize}

\section{The engine beneath the command line}
\label{sec:gr-engine}

\emph{Authority: \path{crates/geolog-chase} (\path{mode.rs},
\path{weight.rs}, \path{view.rs}, \path{delta.rs},
\path{plan_ir.rs}, \path{validity/}), \path{crates/geolog-weights}.
Everything in this section is shipped, tested library surface. The
per-axiom mode partition is the one piece the \texttt{geolog} verb
now reaches — through the surface marks of
\S\ref{sec:gr-chase-modes}; the rest is not reachable from the verb
yet.}

\paragraph{Per-axiom modes, and Enforced.} The engine's mode system
is a per-axiom assignment over the trichotomy \textbf{Chased} /
\textbf{Monitored} / \textbf{Enforced}, with an eligibility-keyed
ambient default (chase-eligible $\Rightarrow$ Chased, everything else
Monitored — reproducing the legacy run-level default, so that default
is semantics-preserving) and validated explicit overrides: a Chased
override on an ineligible axiom is a typed error; Monitored/Enforced
overrides on \emph{eligible} axioms are honoured by a firing
partition — the axiom never fires, and its would-be firings surface
as menu-of-one branch points (the forced repair, surfaced, never
auto-fired). The surface marks lower onto this same machinery, as
a \emph{total} explicit assignment — every axiom's effective mark,
unmarked resolved to \texttt{report} first — so on the file path the
author's word, never eligibility, decides firing, and the ambient
default decides nothing; the ambient default remains the semantics
for library callers that override selectively. Enforced is Monitored
plus the \emph{local write gate}:
an admission verdict over a finished chase result rejects any local
commit whose resulting fixpoint carries an Enforced violation. The
commit-vs-sync seam is the caller's, and its asymmetry is doctrine:
the gate binds \textbf{local commit paths only} — states arriving by
sync or merge are admitted with their judgments recorded, because
gating sync admission on validity would cost true merge and
history-as-authority (\S\ref{sec:gr-sync-merge}). Modes differ in
\emph{when} validity is enforced, never in what validity is. The
adjudicator refuses to grade a mode assignment against a run that was
not partitioned under that same assignment — the silent-vacuity class
(an axiom fired in phase 1 repairing its violations before any
judgment could exist) is rejected loudly, never admitted vacuously.
The meta-schema carries the same trichotomy as rows
(\texttt{Axiom/monitored}, \texttt{Axiom/enforced}, at most one per
axiom; absence of both is the ambient default) —
\S\ref{sec:gr-meta}.

\paragraph{Weights.} The engine is parameterized over an
ordered-semiring provenance-weight algebra
(\path{crates/geolog-weights}): a weighted run integrates each derived
row's weight from its derivation structure, a reweight pass
recomputes after tombstoning, and death-at-zero holds — the
weighted-alive set after tombstone + reweight equals a from-scratch
Boolean chase on the surviving facts (proptested as a differential).
The Boolean rung \emph{is} the unweighted engine: \texttt{run\_chase}
stays pinned to $W = \texttt{bool}$ and is byte-identical to the
weight-blind engine.

\paragraph{Deltas, maintained views, and the plan IR.} The delta bus
makes the engine's unit of change data (element, row, and function
deltas in batches). The view engine compiles violation views —
the failing set of each denial, kept current under row deltas, never
recomputed from scratch — and verifies itself against from-scratch
reference evaluation (a divergence is a reportable defect, and the
rebuild invariant of the repository's battery backs it at replay
scale). Compiled views lower to a tensor-contraction plan artefact
with canonical, digest-citable bytes, a dense reference evaluator,
and an accelerator emission harness; the reference evaluator is the
verdict authority throughout (the lane law,
\S\ref{sec:gr-ground}).

\paragraph{Per-commit validity.} The validity layer
(\path{crates/geolog-chase/validity/}) assembles per-commit validity
judgments over chase results: which commits does this judgment
implicate — well-posed exactly because every derived row and every
invented element carries its firing history
(\S\ref{sec:gr-chase-restricted}). Wiring it into the storage
pipeline end-to-end is design tense (the factoring is committed; the
Tutorial states the same status).

\section{The meta-schema}
\label{sec:gr-meta}

\emph{Authority: \path{theories/GeologMeta.geolog} — the meta-schema's
normative statement, from whose text this section is condensed; where
they disagree, that file wins. The committed genesis golden
(\S\ref{sec:gr-goldens}) is its reference instance; the generator
tool (\path{theories/tools/geologmeta_gen.py}) re-derives and
byte-compares both directions as a drift alarm in CI.}

geolog is a logical framework, not an expressive logic. GeologMeta is
the fixed schema \emph{of} geometric theories: one meta-schema,
frozen at genesis; every theory — including the proof calculi the
companion repositories ship — is an ordinary instance: delta rows
against GeologMeta, carried by ordinary signed commits. Schema,
extension, annotation, and migration therefore live in the same
substrate they govern (the Tutorial's ``no second database'').

\subsection{Sorts and rows}
\label{sec:gr-meta-rows}

Nine schema sorts: \texttt{Theory}, \texttt{Sort}, \texttt{Fun}
(function \emph{or} relation — relationality is a mark, not a second
sort), \texttt{Field}, \texttt{Axiom}, \texttt{Var},
\texttt{Disjunct}, \texttt{Atom}, \texttt{Term}.
Table~\ref{tab:gr-meta-rows} carries the declaration vocabulary over
them, with partiality as documented there (``partial'' in its notes
column means the row may be absent, and absence is meaningful —
Codd's I-mark discipline \cite{codd86}, never ``unknown yet'').

\begin{table}[htbp]
\centering\footnotesize
\begin{tabular}{@{}ll>{\raggedright\arraybackslash}p{0.40\textwidth}@{}}
\toprule
subject & rows & notes \\
\midrule
\texttt{Theory} & \texttt{TheoryExtends[child, parent]} & multi-parent legal; extends graph acyclic \\
 & \texttt{TheoryRevises[new, old]} & revision intent; acyclic; causal half lives in deps \\
 & (no \texttt{Theory/name}) & identity is the element; discovery naming lives above \\
\texttt{Sort} & \texttt{Sort/name}, \texttt{Sort/theory} & \\
 & \texttt{Sort/immutable} & carrier opaque to the chase (choose, never create) \\
 & \texttt{Sort/builtin\_code} & partial: builtin tier only (\S\ref{sec:gr-meta-tier}) \\
\texttt{Fun} & \texttt{Fun/name}, \texttt{Fun/theory} & \\
 & \texttt{Fun/relational} & the relation mark; \texttt{Prop} is not a sort \\
 & \texttt{Fun/codomain} & partial: exactly the non-relational Funs \\
 & \texttt{Fun/arg\_sort} & partial: exactly the unary Funs \\
 & \texttt{Fun/side\_condition} & bounded-computation discharge; entails relational \\
\texttt{Field} & \texttt{Field/fun}, \texttt{Field/name}, \texttt{Field/sort} & name-keyed; canonical order is bytewise name order \\
\texttt{Axiom} & \texttt{Axiom/theory}, \texttt{Axiom/name} & name partial, recommended \\
 & \texttt{Axiom/monitored}, \texttt{Axiom/enforced} & mode rows (\S\ref{sec:gr-engine}); at most one; surface \texttt{report}/\texttt{refuse} lower here, \texttt{chase} to neither \\
\texttt{Var} & \texttt{Var/premise\_of} $\to$ \texttt{Axiom} & universal variables, scoped to the whole axiom \\
 & \texttt{Var/bound\_in} $\to$ \texttt{Disjunct} & existential variables, scoped to one disjunct; exactly one home (partition law) \\
 & \texttt{Var/sort}, \texttt{Var/name} & name is presentation only, hash-excluded \\
\texttt{Disjunct} & \texttt{Disjunct/axiom} $\to$ \texttt{Axiom} & one branch of the conclusion; zero disjuncts on an axiom $=$ rejection rule \\
\texttt{Atom} & \texttt{Atom/premise\_of} $\to$ \texttt{Axiom} & premise atoms hang on the axiom itself \\
 & \texttt{Atom/concl\_of} $\to$ \texttt{Disjunct} & conclusion atoms hang on a disjunct, never on the axiom; exactly one home \\
 & \texttt{Atom/eq\_lhs}, \texttt{Atom/eq\_rhs} & equation atoms \\
 & \texttt{Atom/rel}, \texttt{Atom/rel\_arg} & application atoms; exactly one kind \\
\texttt{Term} & \texttt{Term/var} & variable \\
 & \texttt{Term/app\_fun}, \texttt{Term/app\_arg} & application (arg absent iff nullary) \\
 & \texttt{Term/is\_record}, \texttt{RecordField[rec, field, value]} & records; empty record is the explicit mark \\
 & \texttt{Term/lit\_q,\,lit\_z,\,lit\_nat,\,lit\_b,\,lit\_str,} & one typed literal row per primordial dtype \\
 & \texttt{\ \ lit\_shape,\,lit\_ennrat,\,lit\_timestamp} & \\
\bottomrule
\end{tabular}
\caption{The meta-schema's declaration vocabulary, by subject sort
(\S\ref{sec:gr-meta-rows}).}
\label{tab:gr-meta-rows}
\end{table}

Domain shapes are discriminated by which rows exist —
nullary: neither \texttt{Fun/arg\_sort} nor \texttt{Field} rows;
unary: \texttt{Fun/arg\_sort}; record: \texttt{Field} rows.

The axiom encoding deserves a picture, because it has one level the
sequent notation hides (Figure~\ref{fig:gr-axiom-rows}). The premise
hangs directly on the \texttt{Axiom} element: premise atoms by
\texttt{Atom/premise\_of}, universal variables by
\texttt{Var/premise\_of}. The conclusion sits one level deeper. A
conclusion is a disjunction of existentially quantified conjunctions
(Definition~\ref{def:gr-geometric}), so each branch is its own
\texttt{Disjunct} element, tied to the axiom by
\texttt{Disjunct/axiom}; a conclusion atom names its
\emph{disjunct} as home (\texttt{Atom/concl\_of}), never the axiom,
and an existentially bound variable does the same
(\texttt{Var/bound\_in}). An axiom's premise is therefore the
\emph{set} of atoms homed on it, its conclusion the \emph{set} of its
disjuncts, and its repair arity (\S\ref{sec:gr-chase-classify}) is
its \texttt{Disjunct} count: zero is the empty disjunction, a
rejection rule; one is the canonical, chase-eligible shape; two or
more is a repair menu.

\begin{figure}
\centering
\begin{tikzpicture}[
  elem/.style={draw, rounded corners=1.5pt, inner sep=4pt,
               font=\small\ttfamily},
  rowlab/.style={font=\scriptsize\ttfamily, text=accentmid!75!black,
                 inner sep=1.5pt, fill=white},
  arr/.style={-{Stealth[length=1.8mm]}, semithick, draw=black!70,
              shorten >=2pt, shorten <=2pt}
]
\node[elem] (ax) at (-0.4,0) {ax/ship/manifest : Axiom};
\node[elem] (po) at (-4.7,-1.8) {o : Var};
\node[elem] (pa) at (-1.7,-1.8) {[o:\,o]\,Shippable : Atom};
\node[elem] (d)  at (2.9,-1.8) {d : Disjunct};
\node[elem] (cs) at (0.6,-3.4) {s : Var};
\node[elem] (ca) at (4.4,-3.4) {[o:\,o,\,s:\,s]\,ShippedVia : Atom};
\draw[arr] (po) -- (ax.south west) node[rowlab, pos=0.5, left=4pt] {Var/premise\_of};
\draw[arr] (pa) -- (ax.south) node[rowlab, pos=0.25, right=4pt] {Atom/premise\_of};
\draw[arr] (d)  -- (ax.south east) node[rowlab, pos=0.42, right=4pt] {Disjunct/axiom};
\draw[arr] (cs) -- (d.south west)  node[rowlab, pos=0.5, left=4pt] {Var/bound\_in};
\draw[arr] (ca) -- (d.south east)  node[rowlab, pos=0.5, right=4pt] {Atom/concl\_of};
\end{tikzpicture}
\caption{The Tutorial's \texttt{ax/ship/manifest} ($\forall o.\;
\mathit{Shippable}(o) \vdash \exists s.\; \mathit{ShippedVia}(o,s)$)
as meta-schema rows. Premise structure homes on the \texttt{Axiom}
element; conclusion structure homes on a \texttt{Disjunct}, one
element per branch of the disjunction. Which side of the turnstile
an atom or variable lives on is read off the function row it
carries, and the axiom's repair arity is its \texttt{Disjunct}
count — here one, the canonical shape.}
\label{fig:gr-axiom-rows}
\end{figure}

Field-name uniqueness per
Fun and \texttt{RecordField} uniqueness per (record, field) are load
rejections, deliberately not unification axioms: in a store whose
rows are signed declarations, silently merging independently minted
declarations would erase provenance — collisions are loud.

The meta-axioms, exhaustively (all fifteen; three shapes recur —
entailment, rejection pair, and the $\langle$disjoint,
exhaustive$\rangle$ \emph{partition law} pair stating ``exactly
one''):

\begin{center}
\footnotesize
\begin{tabular}{lll}
\toprule
axiom & shape & states \\
\midrule
\texttt{ax/meta/builtin\_immutable} & entailment & coded sorts are immutable \\
\texttt{ax/meta/side\_condition\_relational} & entailment & side-condition $\Rightarrow$ relational \\
\texttt{ax/meta/relational\_no\_codomain} & rejection & a relation never has a codomain \\
\texttt{ax/meta/codomain\_exhaustive} & disjunctive & every Fun is relational or has a codomain \\
\texttt{ax/meta/arg\_sort\_fields\_disjoint} & rejection & unary and record domains exclude each other \\
\texttt{ax/meta/mode\_marks\_exclusive} & rejection & at most one explicit mode per axiom \\
\texttt{ax/meta/var\_home\_\{disjoint, exhaustive\}} & partition & every Var has exactly one home \\
\texttt{ax/meta/atom\_home\_\{disjoint, exhaustive\}} & partition & every Atom sits on one side of $\vdash$ \\
\texttt{ax/meta/atom\_kind\_\{disjoint, exhaustive\}} & partition & equation or application, never both/neither \\
\texttt{ax/meta/term\_\{var\_app, var\_record, app\_record\}\_disjoint} & rejection & representative term-kind disjointness \\
\bottomrule
\end{tabular}
\end{center}

\noindent (\texttt{ax/meta/codomain\_exhaustive} is also the schema's
own fixture for a disjunctive conclusion with an existential; the
three term-kind rejections are representatives — the full mechanical
pairwise family, the 28 literal pairs included, is elaboration-time
conformance rather than schema text.)

\paragraph{The testable negative.} What is deliberately absent is
structural, not policed: no Formula sort and no
Conj/Disj/Exists/Not/Iff nodes (nesting beyond geometric normal form
is unrepresentable); no negation node and no iff node (both are
authoring-side moves, \S\ref{sec:gr-normalize}); no
\texttt{Index}/dependent-type node (erased upstream); no
\texttt{InfDisj} node (\texttt{Sort/immutable} on the binder's sort
carries the distinction); no proof or derivation sorts (proofs of
object logics are rows of object theories — the logical-framework
doctrine); no weight or provenance slot (engine configuration, not
schema); no annotation vocabulary (annotations are ordinary relations
in ordinary extension theories — metadata needs rows, never a schema
revision).

\subsection{The primordial sort tier}
\label{sec:gr-meta-tier}

Thirteen builtin ground sorts are declared as ordinary Sort rows in
the genesis instance, marked immutable, and fixed forever:
\texttt{Q}, \texttt{Z}, \texttt{Nat}, \texttt{B}, \texttt{Str},
\texttt{ENNRat}, \texttt{Shape}, \texttt{Timestamp}, \texttt{Bytes},
\texttt{PeerId}, \texttt{CommitId}, and the genesis-pinned
extension-band straddle pair \texttt{Date} and \texttt{Interval}.
Beside the thirteen sorts stand three sort \emph{formers}, structural
rather than declared — no \texttt{Sort} row, no declaring element;
like the naturals themselves, they are simply available:

\begin{itemize}
\item \textbf{\texttt{Fin}} — the bounded-index former, an indexed
family $\texttt{Nat} \to \textsf{Sort}$: for each natural $n$,
$\texttt{Fin}(n)$ is the $n$-element index sort. These are the sorts
that index tensor axes (an axis of extent $n$ is indexed by
$\texttt{Fin}(n)$); the surface grammar's bounded-index type
(\texttt{$n$ Index}, \S\ref{sec:gr-decls}) expresses the same idea
at authoring scale, though today the normalizer erases it to
\texttt{Nat} plus range obligations (\S\ref{sec:gr-normalize}).
Cell code \texttt{0x09} (\S\ref{sec:gr-cells}).
\item \textbf{\texttt{Tensor}} — the tensor former,
$\texttt{Shape} \times \textsf{Sort} \to \textsf{Sort}$: given a
shape (a \texttt{Shape} value, the list of axis extents) and an
element sort, the sort of tensors of that shape with entries at that
sort. The element sort may be \emph{any} sort, declared or builtin,
not only a numeric one. Canonical form is flattened, and a rank-zero
tensor collapses to its single entry. Cell code \texttt{0x0A}.
\item \textbf{\texttt{List}} — the finite-sequence former,
$\textsf{Sort} \to \textsf{Sort}$, length carried per value:
$[\![\, x\ \texttt{List} \,]\!] = \sum_{n:\mathbb{N}}
\bigl(\texttt{Fin}(n) \to [\![\, x \,]\!]\bigr)$ — each value brings
its own $n$, the \texttt{DynTensor} mechanism with rank pinned at
one. Positional order is the value (no sorting, no deduplication),
and there is no flatness or unit law: \texttt{(x List) List} is a
distinct, ragged two-level sort, never collapsed. Element scope at
this tier is any sort, nesting included; the one adjudicated
exclusion runs the other way — a list instance is not admissible as
a \texttt{Tensor}/\texttt{DynTensor} \emph{element} in v1, while
\texttt{(x Tensor) List} is legal. Former code \texttt{0x10}; the
surface spelling (\S\ref{sec:gr-decls}) narrows the element scope to
builtin ground sorts; value cells ride dtype code \texttt{0x50}
(\S\ref{sec:gr-cells}).
\end{itemize}

\noindent A fourth structural former stands beside these three at the
wire tier, under its own sort-reference tag rather than a former
code: \textbf{\texttt{DynTensor}} — $\textsf{Sort} \to
\textsf{Sort}$, the \emph{dynamic} tensor former, shape carried in
the value rather than the type (the $\Sigma$-over-shapes reading of
\texttt{Tensor}; tag \texttt{0x36},
\path{crates/geolog-envelope/src/former.rs}). \texttt{Tensor} is the
static reading — shape in the type — and the two are distinct sorts
on purpose, never silently interconverted. Of the tensor pair, only
\texttt{DynTensor} has a surface spelling
(\S\ref{sec:gr-decls}); tensor values are side-condition payload
throughout (\S\ref{sec:gr-ground}).

The tier is ambient: every theory may use it — formers
included — in field and codomain positions without an extends edge.
Minting an element of a builtin sort is a load rejection; builtin
values enter as literals in value cells (\S\ref{sec:gr-cells}).

Genesis-tier ground \emph{vocabulary} equips the primordial carriers
with algebraic structure as ordinary declarations (never as
type-former firmware), each family's evaluation semantics pinned in
\path{theories/GeologMeta.geolog} to the corner case:

\begin{itemize}
\item \textbf{$\mathbb{Q}$ field operations} — \texttt{q\_add}, \texttt{q\_neg},
\texttt{q\_sub}, \texttt{q\_mul}, \texttt{q\_inv}, \texttt{q\_div}:
exact arithmetic on canonical rationals, no rounding anywhere;
partiality is refusal (\texttt{q\_inv(0)} and division by zero refuse
— a field has no answer there, and inventing one is what the exact
tier exists to refuse).
\item \textbf{ENNRat quantale operations} — on the Lawvere rationals
$\mathbb{Q}_{\geq 0} \cup \{\infty\}$: \texttt{ennrat\_add},
\texttt{\_min}, \texttt{\_max}, truncated minus (\texttt{\_monus},
the residual of addition), \emph{both} multiplications
(\texttt{\_mul\_conj} with $0 \cdot \infty = 0$ and
\texttt{\_mul\_disj} with $0 \cdot \infty = \infty$ — the two
conventions declared under distinct names, never silently
harmonized), inversion, residuation, harmonic sum, the two
qualitative-mode side-conditions, and exact comparisons.
\item \textbf{Timestamp} — instants, i64 epoch-microseconds on the
POSIX/UTC scale (timezone is presentation, never value); the numeric
order as side-conditions; checked displacement arithmetic
(\texttt{ts\_add\_us}; overflow refuses, never wraps).
\item \textbf{Date and Interval} — PostgreSQL semantics, as a tested
claim: a corpus of goldens captured from a live PostgreSQL 17
instance is the conformance battery for the calendar family
(\S\ref{sec:gr-ground}); intervals are the three-field
months/days/micros value verbatim, never flattened (comparison uses
the justified collapse, values do not); refusal wherever PostgreSQL
refuses, finite-only.
\item \textbf{Str, CommitId} — disequality only
(\texttt{str\_ne}, \texttt{cid\_ne}): string ordering would smuggle a
collation choice (presentation, like timezone), and ordering commit
hashes would invite misreading bytes as causal order.
\end{itemize}

\noindent Every \emph{other} value format — the floating-point ladder
included — is an extension ground format: declared by an ordinary
signed revision (never genesis, never a schema revision), carrying a
parametric spec over ordinary roster rows (byte width, exponent and
mantissa bits, bias, NaN discipline), with wire codes from the
extension band and operation dispatch bound to the declaring
revision. \texttt{ax/meta/builtin\_immutable} quantifies over every
coded sort, so extension formats are immutable by the same axiom that
binds the genesis roster.

\subsection{Theories as rows; revisions}
\label{sec:gr-meta-revisions}

A theory arrives as a \emph{revision}: a sealed batch of meta-schema
rows in one commit body, signed and content-addressed like every
other commit (\S\ref{sec:gr-commits}), dependencies named. Revisions
publish whole — an axiom and all its structure land within one
revision's commit set, which is what makes zero-disjuncts-as-rejection
well-defined. Identity is structural: a revision hash covers the
declaration content (with presentation-only rows, such as
\texttt{Var/name}, excluded), and a function's signature hash covers
its field names, sorts, and codomain — the identity the ground
seam's dispatch pins (\S\ref{sec:gr-ground}).

Shipped machinery (\path{crates/geolog-meta-rows},
\path{crates/geolog-revision}, \path{crates/geolog-theories}): the
theory$\to$rows compiler (byte-identical round-trips on the committed
corpora); the sealed-revision conformance loader, whose rejections
are loud, named, and collected (one committed negative fixture per
live rule); the structural revision and signature hashes with an
adversarial vector battery; declaration indexes and cross-revision
name resolution; the verified ground-format registry loader; and the
programmatic revision-body builder, whose canonical-form invariants
(backward local references, bytewise field order, exactly-once name
rows) are unrepresentable or checked at build time. The
\texttt{geolog rows} verb (\S\ref{sec:gr-cli}) prints the declaration
slice of this encoding for any theory file; the committed listings
under \path{theories/generated/} carry complete worked theories
(the Petri-net theory is 1{,}239 ops); and \texttt{geolog seal}
(\S\ref{sec:gr-cli}) seals a source file through this machinery —
one revision commit carrying the theories as meta-schema rows (its
dependency table names the genesis and the ground-comparisons
golden, \S\ref{sec:gr-goldens}, uniformly — present whether or not
the source uses the comparison roster, so the table is never
input-dependent), plus a seal companion carrying the source and
registry text — deterministically, so re-sealing an unchanged file
is a no-op (verified: the second \texttt{seal} reports
\texttt{store: genesis already present, comparisons already present,
revision already present, seal already present}).

\section{Commits: envelopes, bodies, cells, storage}
\label{sec:gr-commits}

\emph{Authority: \path{crates/geolog-envelope} (the wire),
\path{crates/geolog-store} (the log),
\path{crates/geolog-ceremony} (the goldens). This section states the
shipped encoding — \emph{encoding v1} — at the grain a consumer
needs; the crate module docs carry the byte-exact layouts, and a
committed cross-language reference implementation backs the golden
vectors.}

\subsection{The skeleton: hash and signature surface}
\label{sec:gr-skeleton}

A commit is a two-layer envelope. The \textbf{skeleton} is the
hash/sign surface: a fixed-layout structure carrying

\begin{center}
\small
\begin{tabular}{ll}
\toprule
field & content \\
\midrule
\texttt{version} & the encoding version byte (1) \\
\texttt{author} & ed25519 public key, 32 bytes \\
\texttt{deps} & the commit-reference list: 32-byte commit hashes \\
\texttt{parent\_count} & how many leading deps are parents \\
\texttt{ts\_claim\_us} & claimed timestamp, epoch-µs (a claim, not a truth) \\
\texttt{message} & optional UTF-8 bytes \\
\texttt{body\_hash} & blake3 of the body bytes \\
\bottomrule
\end{tabular}
\end{center}

\noindent \texttt{deps} is one list in two segments — parents first,
then non-parent content dependencies, each segment sorted strictly
bytewise-ascending, duplicates rejected. \emph{Commit identity is
content:} $\texttt{commit\_hash} = \mathrm{blake3}(\text{skeleton
bytes})$, and the ed25519 signature covers exactly those bytes under
a fixed versioned context prefix
(\texttt{geolog-s4-skeleton-sig-v1}), so no future encoding change
can collide signing domains. Canonicality is definitional, in both
directions: the writer canonicalizes; the reader, after structural
validation, re-encodes the decoded value and requires byte equality
with the input (the fixed-point gate) — one value, one byte string,
independent of any serializer freedom. A signed skeleton can only be
constructed against a body receipt obtainable by durably storing the
body, so the format's sharpest failure mode — a signed skeleton whose
body is absent, a permanently unverifiable dangling commit — cannot
arise from the safe API's normal use.

\subsection{The body: ops}
\label{sec:gr-body}

The \textbf{body} is the ops payload. Wire shape:
$\texttt{u16le body\_format} \parallel \text{frame bytes}$; the
format registry has exactly one live entry, and the numbering below
is the public truth for this repository:

\begin{itemize}
\item \texttt{0} — retired before any production bytes existed (a
migration lane for a legacy dialect; no wire ever carried it).
\item \texttt{1} — \texttt{BODY\_FORMAT\_HEXANE\_TEMPLATES\_V1}:
\textbf{the only format, and the default.} The canonical op sequence
is factored into a template dictionary (one row per cluster of alike
ops) plus columnar hole carriers; everything the dictionary hoists is
spelled once per cluster, and each op row pays only for its hole
values. The factoring is fully \emph{derived} from the content —
clustering, per-position carrier modes, and traversal order are laws,
not producer choices — and a frame spelling the same rows any other
way is refused at decode. The frame ends with the same fixed-point
gate the skeleton uses: decode re-encodes and requires byte equality.
$\texttt{body\_hash} = \mathrm{blake3}(\text{body bytes})$.
\end{itemize}

\noindent The op vocabulary:

\begin{center}
\small
\begin{tabular}{lll}
\toprule
op & payload & meaning \\
\midrule
\texttt{elta} & sort reference & mint an element of a (mutable) sort \\
\texttt{rela} & relation ref, argument cells & assert a relation row \\
\texttt{funa} & function ref, argument cells, value cell & assert a function value \\
\texttt{eltd} & target reference & delete an element (never a meta-row target) \\
\texttt{blob} & format tag, opaque bytes & content-addressed bulk (\S\ref{sec:gr-cells}) \\
\texttt{unknown} & kind tag, opaque bytes & forward-verifiability carry \\
\bottomrule
\end{tabular}
\end{center}

\noindent References are either \emph{local} (an earlier op of the
same body, backward only) or \emph{external} (content-addressed,
through the deps); names are data (\texttt{Sort/name} rows), never
identity. An unknown op kind is preserved byte-for-byte and carried,
never dropped and never an error — a reader that predates a kind
still parses correctly, because no decode path may silently drop
data. Decode enforces hostile-input hygiene throughout: claimed
counts are bounded (the op-count ceiling is $2^{20}$ per body) and
bounds-checked before any allocation.

\subsection{Cells and the dtype registry}
\label{sec:gr-cells}

Argument and value positions carry \emph{cells}: element references,
or literals of the primordial dtypes. The primordial band is complete
— \texttt{0x00}–\texttt{0x0F}, no holes:

\begin{center}
\small
\begin{tabular}{llll}
\toprule
code & dtype & code & dtype \\
\midrule
0x00 & \texttt{Q} & 0x08 & \texttt{Shape} \\
0x01 & \texttt{Z} & 0x09 & \texttt{Fin} (former) \\
0x02 & \texttt{Nat} & 0x0A & \texttt{Tensor} (former) \\
0x03 & \texttt{B} & 0x0B & \texttt{Timestamp} \\
0x04 & \texttt{Str} & 0x0C & \texttt{Bytes} \\
0x05 & \texttt{ENNRat} & 0x0D & \texttt{BlobRef} \\
0x06 & \texttt{ZBig} & 0x0E & \texttt{PeerId} \\
0x07 & \texttt{NatBig} & 0x0F & \texttt{CommitId} \\
\bottomrule
\end{tabular}
\end{center}

\noindent \texttt{ZBig}/\texttt{NatBig} are multiprecision
realizations of $\mathbb{Z}$/$\mathbb{N}$ (distinct codes because
widening a code's grammar would be reinterpretation, which the
registry's general law forbids); \texttt{BlobRef} is the out-of-line
spelling of \texttt{Bytes}/\texttt{Str} above the inline threshold (a
35-byte content-address cell); \texttt{CommitId} is a reference — a
tiny index into the enclosing skeleton's deps list — and never a
context-free literal. The extension band \texttt{0x10}–\texttt{0x3F}
carries ground formats (the float ladder \texttt{0x10}–\texttt{0x16};
the genesis-pinned \texttt{Date} \texttt{0x17} and \texttt{Interval}
\texttt{0x18}), decoded opaquely and validated at closure tier; the
absence band \texttt{0x40}–\texttt{0x4F} carries typed absence.
Values of the list former ride dtype code \texttt{0x50} — the first
assignment past the reserved bands, sitting out of numeric order
because the primordial band was already complete and the former
registry's \texttt{0x10} (\S\ref{sec:gr-meta-tier}) falls inside the
extension band — validated at decode like the primordials, in two
first-byte-discriminated arms: an \emph{inline} arm (element dtype,
count, then length-framed canonical element payloads — scalar
elements only: the primordial literal codes minus \texttt{Tensor},
\texttt{BlobRef}, and \texttt{CommitId}) and an \emph{envelope} arm
(a 34-byte multihash content address under a dedicated derive-key
context), partitioned at 64\,KiB in the tensor family's style;
nested lists and non-scalar elements ride the envelope arm, and any
cell code $\geq$ \texttt{0x51} is a decode error.
Over-threshold bulk travels as \texttt{blob} ops (format 0 =
\texttt{blob-raw-v1}) on the same signed-envelope machinery —
signing, deps, sync, and equivocation detection inherited, interior
deliberately opaque.

\subsection{The store log}
\label{sec:gr-store}

The persistent store (\path{crates/geolog-store}) appends envelopes
and bodies to a log file with a fixed grammar:

\begin{Verbatim}[fontsize=\small,frame=leftline,rulecolor=\color{accentmid},framesep=3mm,xleftmargin=3mm]
file    := header record*
header  := b"glog2607" | u32-LE format_version (1) | u32-LE reserved (0)
record  := u32-LE payload_len | u8 kind | payload
kind 1  := ENVELOPE: sig[64] | skeleton bytes
kind 2  := BODY:     raw body bytes (identity = blake3(bytes))
\end{Verbatim}

\noindent There is no per-record checksum, by adjudicated design:
every envelope is signature-verifiable and every body
content-hash-verifiable — stronger than any CRC — and a checksum
cannot repair length-prefix desynchronization anyway. A record length
above the 64\,MiB ceiling is treated as torn or garbage bytes misread
as a length, not as a record. The grammar layer never classifies a
failure as crash artefact versus corruption — that is recovery
policy, one layer up, which consults a committed hint the grammar
deliberately knows nothing about. Header deviations (magic, version,
nonzero reserved field) refuse loudly; a future format is a new
version behind that gate, never a silent reinterpretation.

\subsection{The committed goldens and the ceremony}
\label{sec:gr-goldens}

Four signed golden envelopes ship under \path{theories/}: the
\textbf{genesis} (GeologMeta's own declaration rows as a signed
commit, the root of the graph, $\texttt{deps} = []$, message
\texttt{geolog-meta-genesis-v1}); the \textbf{ground-formats}
revision (the float-format vocabulary, $\texttt{deps} =
[\text{genesis}]$); the \textbf{zones} vocabulary
($\texttt{deps} = [\text{genesis}]$); and the
\textbf{ground-comparisons} roster ($\texttt{deps} =
[\text{genesis}]$, message \texttt{geolog-ground-comparisons-v1}) —
the builtin ground-comparison vocabulary (\texttt{q\_lt},
\texttt{q\_le}, \texttt{q\_ne}, \texttt{nat\_lt}, \texttt{nat\_ne})
as a sealed revision, the second dependency of every revision commit
\texttt{geolog seal} mints (\S\ref{sec:gr-meta-revisions}). Each is a deterministic
function of its committed structure-normative listing and a
\emph{public} seed committed in source: anyone can derive the signing
key and reproduce every byte, and
\texttt{cargo run -q -p geolog-ceremony -- verify --root theories}
does exactly that in every CI pass. The key certifies nothing about
who signed; what makes these goldens canonical is that they are
byte-reproducible from the committed listings by the committed tool.
A real deployment mints its own genesis under its own key; the
ceremony crate is the complete worked example of how. The genesis is
also the meta-schema's homoiconic fixed point — GeologMeta described
as GeologMeta rows — and the generator's \texttt{check} mode
re-derives it from \path{theories/GeologMeta.geolog} both ways as a
drift alarm.

\section{Sync, merge, and equivocation}
\label{sec:gr-sync}

\emph{Authority: \path{crates/geolog-sync} (whose README freezes the
contract: objects, invariants, the frontier exchange, the convergence
theorem, and a ledger of what is and is not yet claimed) and
\path{crates/geolog-envelope/src/equivocation.rs}.}

\subsection{Merge}
\label{sec:gr-sync-merge}

\begin{definition}[merge]
\label{def:gr-merge}
The state of a replica is a set of signed commits (a deps-closed
commit bag). Merging replicas is set union on commits followed by
re-derivation of all derived state. Merge is therefore total,
commutative, associative, and idempotent; it never conflicts
structurally, and admission order does not affect derived state.
\end{definition}

\noindent A merged state may satisfy some axiom's violation premise —
that is a judgment (\S\ref{sec:gr-chase-judgments}), repaired by
further ordinary commits, never by a conflict-resolution protocol.
The Enforced gate binds local commit paths only
(\S\ref{sec:gr-engine}); gating sync admission on validity is the
classic failure this architecture refuses, because it would cost true
merge and history-as-authority. Identity questions inside the merged
state (two writers' records denoting one thing) are the theory's
business, decided by its equality axioms and computed by congruence
(\S\ref{sec:gr-chase-congruence}). The rebuild invariant backs the
order-independence claim mechanically: the repository's battery diffs
incrementally maintained state against from-scratch replay, and a
divergence is a release-blocking defect.

\subsection{Replication}
\label{sec:gr-sync-replication}

Two shipped strata. \textbf{Sync v0 — the file mesh}: whole-file,
N-device mesh sync of per-device, single-writer, append-only envelope
logs plus content-addressed body directories; per-source mirrors and
per-run receipts (``silence means broken''); a torn trailing frame is
classified by comparing against the prior round's observation
(growing source: benign; stable source with a persisting tear: a
problem) — a distinction a single snapshot cannot make.
\textbf{Replay-through-ingestion} is the only path synced envelopes
enter a local view: re-verification through the envelope crate's own
verify path (no new trust surface), idempotent by commit hash,
order-free via a deps-frontier worklist fixpoint.

\textbf{Sync v1 — the connected protocol}, layered strictly above v0:
a deterministic integer-keyed message codec (strict known-kind
decode, unknown-kind skip), frontier-exchange pulls, subscribe push,
acknowledged head updates, credit-based flow control, bounded
staging, and one door: admission through the same replay path.
Equivocation detection runs incrementally during sync and is loud,
never a gate.

\paragraph{Realms and partitions.} Sync is scoped by \emph{realm} —
the wire's opaque slice key, today a plain partition name: sessions
pull and subscribe per realm, and a realm's head set is an exact
summary of its deps-closed store (the crate's \textsc{inv-closure},
which is store-global). A materialized partition must be
dependency-closed and revision-closed (no partition splits a semantic
revision unit); the reference builder and its checker are separate
artefacts, and the checker re-verifies against the store rather than
trusting the builder. Dependencies that cross a realm boundary are
governed by a per-realm \texttt{TransportMode} dial:
\texttt{refused} (the default), \texttt{reported} (the crossing
becomes a declared boundary reference), or \texttt{repaired} (a
locally held foreign dependency's pinned closure joins; a repair that
would need a network fetch is a typed refusal). The \texttt{repaired}
dial also binds at admission: a realm registered \texttt{repaired}
admits entries whose dependencies are admitted under \emph{other}
realms — the dependencies must still be admitted locally, so closure
remains store-global — and the dial is local policy, taken at store
open and never negotiated on the wire; an undialed peer stages the
same entries, a liveness asymmetry rather than a safety one.

\subsection{Equivocation}
\label{sec:gr-sync-equivocation}

\begin{definition}[equivocation-suspicious pair]
\label{def:gr-equivocation}
Two distinct envelopes from one author, neither of which causally
references the other (through \texttt{deps}), form an
equivocation-suspicious pair.
\end{definition}

\noindent Detection is structural — from the bytes alone, no
consensus protocol, no quorum — and local to whatever envelope set is
supplied: a dep pointing outside the set is ``unknown, not
reachable,'' never an error. Two different-bytes envelopes claiming
one commit hash refuse loudly (a genuine hash collision or a hostile
input; never silently resolved by picking a winner). The envelope
crate ships the demonstration as a test — two signed commits claiming
the same chain position, exhibited and caught — and the Tutorial
locates this in the systems landscape: the offence a broker cannot
even state is detected here from the data model. The mechanism,
signed hash-linked histories making forks structurally visible, is
the same one that makes CRDTs Byzantine fault tolerant
\cite{kleppmann22}.

\section{The ground seam}
\label{sec:gr-ground}

\emph{Authority: \path{crates/geolog-ground-api} (the contract) and
\path{crates/geolog-ground} (the reference floor, its first
provider).}

Ground vocabulary — relations and functions over immutable carriers,
decided by procedures rather than stored rows — crosses one seam.
\textbf{Declaration is data}: ground vocabulary is declared where all
vocabulary is declared, as committed theory rows, and resolves to a
pinned identity (declaring commit + row index) with a verified
signature and, where one exists, the commit hash of the committed
evaluation-semantics row. \textbf{Implementation is code}: a
\emph{provider} is one value carrying curated claims plus the
evaluator over its own entries. \textbf{The binder is where they
meet}, and it is engine-owned law no provider can bypass or see:
every resolved declaration is adjudicated against every provider's
claims, with three outcomes per entry —

\begin{center}
\small
\begin{tabular}{ll}
\toprule
status & meaning \\
\midrule
Implemented & exactly one claim verified, committed semantics covered — a dense engine id is allocated \\
Deferred & declared-but-unbound, safe-fail: no claimant, no semantics row, or a stale provider \\
Refused & code and theory disagree (signature mismatch) — the drift alarm, distinct from Deferred \\
\bottomrule
\end{tabular}
\end{center}

\noindent Two verified claims on one identity are always an error;
provider order is never semantics. Dispatch binds to pinned
identities, never names: a same-named impostor declaration gets no
oracle and its atoms defer forever. An unbound entry never obtains a
dense id, so a compiled plan citing it cannot even be expressed —
\emph{compile refusal by construction}, the seam's realization of
refusal-over-guessing.

The laws, stated once here and cited by everything that consumes the
seam: \textbf{purity} (every evaluator is stateless and deterministic
per its declared scope — same cells, same outcome, every process,
every architecture; a family whose reference uses floating point must
pin bit-exact step semantics or cannot be admitted);
\textbf{the lane law} (an accelerated lane may change cost, never
verdicts, and ships with a conformance battery against the always-on
scalar reference — a lane that cannot pass its battery does not
dispatch); \textbf{refusal} (refusal is data, and conservative: under
an Enforced axiom, a gate that cannot evaluate must not admit; no
seam method panics on validated cells); \textbf{no enumeration} (a
ground relation filters, never generates; a ground function inverts
only where the algebra genuinely determines the missing argument —
underdetermined means the engine must not enumerate); and
\textbf{scope} (pure is the norm; the one typed exception is
snapshot-scoped structural relations reading instance rows through an
engine-supplied view, with the routing engine-enforced in both
directions). Cells cross the seam in their canonical encodings,
opaque to the contract; there is no intermediate certificate file
format anywhere — a checker is a ground relation evaluating over rows
the engine already speaks.

The \textbf{reference floor} ships as the seam's first provider:
exact arbitrary-precision rational arithmetic, the
extended-nonnegative-rational quantale operations, calendar
arithmetic with documented, PostgreSQL-conformant semantics (the
golden corpus of \S\ref{sec:gr-meta-tier}), and the curated
ground-comparison family. The floor's reference path is
unconditional; the calendar \emph{functions} sit behind a feature
flag whose absence makes their claims absent — declared entries then
defer at bind, so \textbf{verdicts never change with a feature flag}.
Application side-condition families (the proof-calculus checkers of
the companion repositories) are application content: they ship as a
provider crate owned by the application, composed at build time
through this same binder — the only sanctioned coupling between this
repository and its first big application. Reaching the seam from the
\texttt{geolog} command line is design tense; the engine-level seam,
the floor, and their batteries are shipped and tested today (the
Tutorial's ground chapter walks the same line).

\section{Capabilities}
\label{sec:gr-caps}

\emph{Authority: \path{crates/geolog-caps}.}

Authorization arrives as an algebra over commits, not an access
list beside them. Shipped today, as data and replay semantics:
capability operations (delegations and revocations) travel as
ordinary signed revisions; a normative replay computes which revocations are
effective; authorization judgments are query-time and
position-relative. The committed design decisions, decision by decision (the authority
crate's own phrase):

\begin{itemize}
\item a grade lattice
$\text{Define} < \text{Retract} < \text{Admin}$ per sort — element
deletion is the non-monotone direction, so it is strictly higher
privilege than minting: the add/delete split of the op vocabulary,
lifted;
\item a per-relation \texttt{Assert} chain sharing no element with
the sort chain above bottom, so sort authority and row authority can
never satisfy one another;
\item dual Admins per theory and per instance, with theory-Read as
the base prerequisite to holding anything effective on a theory's
instances;
\item a conservative cross-tier stance: no implication in either
direction; dormant grants are first-class, queryable, and surfaced by
an authoring lint that replay never reads;
\item revocations executed in earliest-add seniority (ascending
authority-provenance depth, digest-tiebroken), with the auth-order
projection carried as data so seniority cannot be padded by
packaging;
\item edge identity as (revision, index), so identical re-grants in
distinct revisions are distinct edges.
\end{itemize}

\noindent \emph{Design
tense:} enforcement at the store boundary — wiring these judgments
into admission of local commits — is the committed direction; today
the algebra, replay, and judgments ship as a library with their
property tests, and no store consults them yet.

\section{The \texttt{geolog} command line}
\label{sec:gr-cli}

\emph{Authority: \path{crates/geolog} — \path{main.rs} for the
surface, the modules cited throughout this chapter for each verb's
semantics. The end-to-end tests in \path{crates/geolog/tests/cli.rs}
run the book's own transcripts against the committed examples.}

\begin{Verbatim}[fontsize=\small,frame=leftline,rulecolor=\color{accentmid},framesep=3mm,xleftmargin=3mm]
geolog check     <file.geolog> [--strict]
geolog check     --store DIR [<commit>] [--strict]
geolog chase     <file.geolog> [--instance A[,B...]]
                               [--budget N] [--mode chased|monitored]
geolog chase     --store DIR [<commit>] [--instance A[,B...]] [--budget N]
geolog seal      <file.geolog> --store DIR [--strict] [--repro]
geolog extract   <commit> --store DIR [--referencing] [-o FILE.glc]
geolog restrict  <file.glc> --store DIR --theory SEL [--revision HEX]
                               [-o FILE.glc]
geolog load      <file.geolog> [--strict] [--dry-run] [-d]
geolog load      <file.glc> [--store DIR]
geolog fmt       <file>
geolog rows      <file.geolog>
geolog normalize <file.geolog> [--strict]
geolog census    <file.geolog>... [--rows LISTING]...
                               [--budget N] [--tsv FILE]
\end{Verbatim}

\paragraph{The specified verbs.}
\textbf{check} runs parse $\to$ normalize $\to$ lint on the theory
text and validates every instance block against the elaborated
closure by actually seeding it — exactly the checks \texttt{chase}
would run, without the chase — reporting each stage and then three
summaries: a \emph{mode census}, what the file's axioms will do when
chased (\texttt{modes: 2 chase · 1 report (0 by the unmarked default)
· 0 refuse} for the Quick Start theory — so an all-\texttt{report}
file, deriving nothing by design or awaiting its
explicit-\texttt{chase} port, is legible at a glance); the move
census (\S\ref{sec:gr-normalize}); and, where the chase bridge
carries the fragment, a chase-termination advisory (weak acyclicity)
computed under the file's own surface modes — only the
\texttt{chase}-marked fragment generates, so a theory of
\texttt{report}-marked constraints is termination-trivial by
construction, however generative its axioms would be if fired. With
\texttt{--store}, the same strata run over a sealed source instead
of a file, after signature verification and a coherence gate: the
sealed source must re-derive the sealed revision rows
byte-identically.
\textbf{chase} seeds the selected instance blocks and runs the chase
to a fixpoint, then renders the report of \S\ref{sec:gr-report}.
\textbf{rows} prints the theory's declaration slice as meta-schema
rows in the committed listings' notation (local refs \texttt{L$n$},
genesis identities \texttt{G<…>}), with a closing note counting the
axioms that continue as rows beyond the slice; it requires the
one-field record spelling for unary relations, because fields are
rows with names in the meta-schema (verified: \texttt{relation
`Marked` has a unary (anonymous) domain — … rows emission needs the
one-field record spelling: `Marked : \{x: A\} -> …`}).

\paragraph{The remaining verbs, at citation grain.} Each is shipped
and tested; this manual states its role and leaves its full
semantics to the named authority. \textbf{seal} runs parse $\to$
normalize $\to$ canonical revision rows $\to$ a signed envelope pair
in the store: a revision commit and a seal companion
(\S\ref{sec:gr-meta-revisions}); re-sealing an unchanged file is a
no-op. With \texttt{--repro} the signature uses the public
reproduction seed, making the commit hashes store-independent (pin
reproducibility); such commits carry no authorship and must never
anchor authority. \textbf{extract} writes a commit and its transitive deps from
the store to a \path{.glc} container (magic \texttt{geologfl},
self-contained, verified on read rather than trusted);
\texttt{--referencing} additionally sweeps every store commit whose
deps reach the target, so all commits referencing an instance travel
together. \textbf{restrict} is theory-fragment restriction of a
certificate container — the get side of the delta lens a
\texttt{theory B extends A} induces: it keeps exactly the rows whose
species live in the target theory's extends-closure (derived from
the sealed revision's own declaration rows), projects the
prerequisite table onto the surviving boundary, copies unchanged
commits byte-verbatim, re-seals changed ones deterministically with
the public reproduction key, and refuses — with a named locus — any
container whose kept rows reference stripped content; a
pure-fragment container round-trips byte-identical. \textbf{load}
with a \path{.glc} argument decodes and
verifies the whole container — signatures, body hashes, topological
order — then admits it into \texttt{--store}; with a
\path{.geolog} argument it stages the theory's aggregation surface
into a live engine (\S\ref{sec:gr-editions}). \textbf{fmt} parses a
file, legacy spellings included, and prints the canonical source
respelling — the corpus-port instrument (\S\ref{sec:gr-decls}).
\textbf{normalize} prints the geometric normal form the engine
consumes (\S\ref{sec:gr-gnf}) — the committed-golden regeneration
verb. \textbf{census} is the merge-reliance census over authored
equality heads; its method and honesty ledger live in
\path{crates/geolog/src/census.rs}. Authority for all of these:
\path{crates/geolog/src/main.rs} and the modules it names.

\paragraph{Options.} Classification travels in the source file
(\S\ref{sec:gr-registry}), so the file-mode verbs take the file and
nothing else; \texttt{--registry FILE}, the deprecated migration-era
carrier, refuses beside import lines when its registry carries
classification (a pattern-only registry composes, its aux names
kept clear of imported ones — \S\ref{sec:gr-registry}), and in
store mode the flag is refused outright: the registry text travels
sealed.
\texttt{--strict} rejects the deprecated spellings — bracket
records, binders-before-name theory headers, \texttt{=~\{} instance
bodies, \texttt{:~Sort} declarations — each with a pointer at the
ratified form (\S\ref{sec:gr-decls}) — and demands the
\texttt{import std} header for builtin ground vocabulary used by
name (\S\ref{sec:gr-registry}). \texttt{--instance A[,B…]}
selects blocks by name — the default is every block in the file, as
their union (\S\ref{sec:gr-instances}); naming a missing block is a
usage error that lists the blocks present. \texttt{--budget N} bounds
chase iterations (default 1000); exhaustion is a reported judgment,
not an error. By default \texttt{chase} runs each axiom under its
authored surface mode (\S\ref{sec:gr-chase-modes});
\texttt{--mode chased|monitored} is the legacy run-level flattener
for files whose axioms carry no marks, and refuses — naming the
marked axioms — on any file that does. \texttt{chase} on a file with
no instance block is a usage error that says what to add and points
at \texttt{check} for theory-only validation.

\paragraph{Exit codes.} \texttt{0} — clean: the verb completed and,
for \texttt{chase}, the fixpoint carries no judgments. \texttt{1} —
judgments at fixpoint: the run completed and reported violated
denials, pending branch choices, or an exhausted budget. \texttt{2} —
the input never ran: usage errors, unreadable files, pipeline
rejections, bridge refusals, seeding contradictions, and engine load
refusals. The 1/2 boundary is the error type's own doctrine:
judgments from a completed chase are results, not errors.

\paragraph{Constituent binaries.} The workspace also ships
development CLIs for constituent crates — \texttt{glnorm} (the
normalizer pipeline standalone), \texttt{glstore} (the store log),
\texttt{glsync} (the file-mesh sync) — documented in their crates;
\texttt{geolog} is the carryable verb and the only surface this
chapter specifies.

\section{Editions: what this manual does not yet state}
\label{sec:gr-editions}

Open seams, each named with where its truth lives today; a future
edition of this chapter absorbs them in this order.

\begin{itemize}
\item \textbf{The aggregation surface.} The ratified authoring
grammar beyond the fragment this manual states — \texttt{law}
members (named or anonymous) whose bodies carry contraction index
groups and string-literal contraction specs, the arity-typed
contraction resolver with its refusal roster
(\texttt{E-RANGE}, \texttt{E-SPEC-*}, \texttt{E-RPN-*},
\texttt{E-LAW-*}), and the free-construction instance declarations
(parameters, \texttt{extends} targets, and body rows as three seed
sources) — is shipped in \path{crates/geolog-normalize} (the parser
module's documentation is the compact normative summary) and staged
into a live engine by \texttt{geolog load}, whose emission header
marks the engine realization design-tier pending the
aggregation-schema promotion. This edition states the surface's
lexical floor and type-grammar hooks (\S\ref{sec:gr-lexical},
\S\ref{sec:gr-decls}) and defers the rest to the next.
\item \textbf{The store-backed surface at this manual's grain.}
\texttt{seal}, \texttt{extract}, \texttt{load} of a \path{.glc}
container, the store modes of \texttt{check} and \texttt{chase}, and
\texttt{census} all ship today (\S\ref{sec:gr-cli}); the container
grammar, the store-mode conformance rules, and the census method are
normative in the facade's modules. A future edition states them at
the grain \S\ref{sec:gr-commits} gives the envelope.
\item \textbf{The byte-exact wire grammar.} \S\ref{sec:gr-commits}
states the encoding at consumer grain; the op-payload grammar, the
template-frame channel layout, and the conformance-rule table
(load-rejection classes and their fixtures) are normative today in
\path{crates/geolog-envelope}'s module documentation and
\path{crates/geolog-meta-rows}' loader and negative fixtures.
\item \textbf{Per-commit validity, end to end.} The judgment
assembly ships (\S\ref{sec:gr-engine}); the storage-pipeline wiring
and its report vocabulary are design tense.
\item \textbf{Capability enforcement} at the store boundary
(\S\ref{sec:gr-caps}).
\item \textbf{The ground seam from the command line}
(\S\ref{sec:gr-ground}): the chase-path elaboration of ground sorts
remains refused (\S\ref{sec:gr-elab}), and the \texttt{load} path
stages its ground-vocabulary dependencies as roster rows while its
engine realization stays design-tier (previous bullet).
\item \textbf{Cross-revision \texttt{extends} in the facade}, native
bounded-index types, and the full evaluation-semantics tables of the
genesis ground families (today: pinned, with their corner cases and
golden corpora, in \path{theories/GeologMeta.geolog} and the ground
crates' tests).
\end{itemize}

%   !=@A; ghost element; List field arm; Str List List; List element
